\documentclass[UTF8,a4paper,11pt,fontset=fandol]{ctexart}
\usepackage[margin=2.2cm]{geometry}
\usepackage{amsmath,amssymb,booktabs}
\setlength{\parskip}{0.3em}
\linespread{1.1}

\begin{document}
\pagestyle{plain}

\begin{center}
  {\Large\bfseries 第七题：SOR 迭代与 Cholesky 分解比较}
\end{center}

\section{原题抄录}

用 SOR 迭代法求解方程组
\[
\begin{pmatrix}
4&-1&0\\
-1&4&-1\\
0&-1&4
\end{pmatrix}
\begin{pmatrix}x_1\\x_2\\x_3\end{pmatrix}
=\begin{pmatrix}1\\4\\-3\end{pmatrix}.
\]
其精确解为 $x^*=(0.5,1,-0.5)^{\mathrm T}$。取初值向量
$x_0=(0,0,0)^{\mathrm T}$，分别取 $\omega=1.03$、$\omega=1$、
$\omega=1.1$ 进行计算，要求
\[
\|x_k-x^*\|_\infty\leq\frac12\times10^{-5},
\]
并确定每一个 $\omega$ 值对应的迭代次数。
和 Cholesky 分解直接法求解进行比较。

\section{代码的结构、算法说明}

程序首先定义系数矩阵 $A$、右端向量 $b$ 和教材中的分裂矩阵
$D,L,U$，满足 $A=D-L-U$。其中 $D$ 是对角部分，$L,U$ 分别是
严格下、上三角部分的负值。程序通过
\texttt{numpy.linalg.solve(A, b)} 计算参考解 \texttt{x\_exact}，
供停止条件使用；本题参考解与题目给出的 $x^*$ 一致到浮点数精度。

\texttt{distance\_infinity(x, y)} 返回两个向量之差的无穷范数，
即各分量差的绝对值的最大值。

\texttt{SOR(omega)} 从零向量出发，令
\[
M=D-\omega L,\qquad N=(1-\omega)D+\omega U.
\]
这里是对 $\omega A$ 的分裂，即 $\omega A=M-N$。迭代格式为
\[
x_{k+1}=M^{-1}Nx_k+\omega M^{-1}b
=(D-\omega L)^{-1}\bigl((1-\omega)D+\omega U\bigr)x_k
+\omega(D-\omega L)^{-1}b.
\]
代码用 \texttt{numpy.linalg.solve} 求解相应线性方程组，
而不显式计算 $M^{-1}$。每次更新后判断误差是否小于
$5\times10^{-6}$，满足条件就返回近似解及迭代次数。
初值对应 $k=0$，循环计数从 1 开始，最多进行 99 次迭代。
代码使用严格不等式 $<$，比原题的 $\leq$ 略严格，
但本题三组结果均不处于等号边界，因此不影响达标次数。

\texttt{cholesky(A)} 按下三角形式计算 Cholesky 分解。
该函数内部的局部变量 \texttt{L} 是 Cholesky 因子，与 SOR 的分裂矩阵
$L$ 含义不同；这里记为 $L_{\mathrm C}=(\ell_{ij})$，满足
$A=L_{\mathrm C}L_{\mathrm C}^{\mathrm T}$。计算公式为
\[
\ell_{ii}=\sqrt{a_{ii}-\sum_{s=1}^{i-1}\ell_{is}^{\,2}},\qquad
\ell_{ij}=\frac{a_{ij}-\sum_{s=1}^{j-1}\ell_{is}\ell_{js}}{\ell_{jj}}
\quad(i>j).
\]
\texttt{cholesky\_solve(A, b)} 得到因子后，依次调用
\texttt{numpy.linalg.solve} 求解
\[
L_{\mathrm C}y=b,\qquad L_{\mathrm C}^{\mathrm T}x=y,
\]
从而得到原方程组的数值解。最后，主程序分别调用三组 SOR 和
Cholesky 求解函数，\texttt{print\_vector} 将各数值解统一输出到
小数点后十二位。

\section{输出结果呈现}

运行所给代码，得到以下结果。表中的向量分量保持程序的输出精度；
直接法不计 SOR 迭代次数，以 ``---'' 表示。

\begin{center}
\small
\setlength{\tabcolsep}{4pt}
\begin{tabular}{lccrrr}
\toprule
方法 & $\omega$ & 迭代次数 & $x_1$ & $x_2$ & $x_3$\\
\midrule
参考解 & --- & --- & 0.5 & 1 & $-0.5$\\
SOR & 1.03 & 5
& 0.500004471768 & 1.000001611252 & $-0.499999737298$\\
SOR & 1 & 6
& 0.500003814697 & 1.000001907349 & $-0.499999523163$\\
SOR & 1.1 & 6
& 0.500003630405 & 0.999998540537 & $-0.500000039393$\\
Cholesky & --- & ---
& 0.500000000000 & 1.000000000000 & $-0.500000000000$\\
\bottomrule
\end{tabular}
\end{center}

\section{对输出结果的分析和结论}

三个松弛因子下的 SOR 结果均满足题目要求。$\omega=1.03$ 需要
5 次迭代，$\omega=1$ 和 $\omega=1.1$ 均需要 6 次迭代。
因此，在本题给定的初值、误差限和三个参数中，$\omega=1.03$
所需迭代次数最少。$\omega=1$ 时 SOR 就是 Gauss--Seidel 迭代；
将松弛因子增大到 $1.1$ 没有进一步减少次数，说明松弛因子的选择
需要结合具体问题。

Cholesky 输出在小数点后十二位与精确解一致，数值精度高于
按本题误差限停止的 SOR 结果。这里的显示一致不意味着浮点计算
完全没有舍入误差。由于 $A$ 对称且顺序主子式 $4,15,56$ 均为正，
矩阵对称正定，适合使用 Cholesky 分解。

对于本题的小规模对称正定方程组，Cholesky 法无需选择初值和
松弛因子，通过一次分解及两个三角方程组求解即可得到高精度结果；
SOR 则通过逐步逼近达到指定误差要求。本次输出验证了两种方法
均能有效求解该方程组。程序没有输出运行时间，因此本实验仅比较
数值结果和 SOR 迭代次数，不据此判断两种方法的实际耗时。

\end{document}